\section{讨论}

\subsection{连续正则化优于二值门控}

本文一个关键发现是二值激励门控在反馈控制AUV中不可靠。无论使用$\lambda_{\min}$还是偏航率阈值，控制器持续产生的航向修正都阻止了门控关闭。连续GM正则化从根本上解决了这一问题：通过始终施加弱衰减，它自然地在高激励阶段被梯度更新主导，在低激励阶段发挥正则化作用，无需阈值调谐。

\subsection{精度-物理可接受性权衡}

实验结果揭示了海流观测器评估中的一个重要权衡：RMSE最低的估计并非总是最好的。在高噪声条件下，无限幅变体实现了最低的RMSE（0.126 vs.\ 0.243），但产生了物理不可能的单步跳变（2.9 m/s）。这强调了对海流观测器进行多目标评估的必要性——单步更新量、过冲和最终偏差应作为RMSE的补充指标。

\subsection{噪声驱动与失配驱动的限幅需求}

压力矩阵实验产生了一个关键洞见：限幅需求仅由DVL噪声驱动，而非模型失配。在所有失配水平下，低噪声条件的限幅倍数仅为1.6--1.8倍，而高噪声条件为29--33倍。这一发现挑战了限幅主要用于模型失配鲁棒性的直观假设，转而将其定位为噪声鲁棒性机制。对于模型失配，连续GM正则化（而非限幅）提供了主要保护。

\subsection{安全性视角}

从安全关键AUV运行的角度，PC-RCO可被视为海流观测的"安全包装器"。它不承诺最低的RMSE，但保证：（1）每步估计变化有界；（2）低激励阶段估计向先验回归而非漂移；（3）任何条件下不会出现物理不可能的瞬态估计。在AUV安全运行优先于估计精度指标的应用场景中，这种有界行为具有价值。

\subsection{局限性}

当前研究存在若干局限。首先，验证基于仿真，需要水池试验或海试来确认实际性能。其次，自适应限幅依赖准确的在线噪声估计，在非平稳噪声条件（如DVL底跟踪丢失）下可能退化。第三，方法假设时不变的CFD先验；基于累积运行数据自适应更新先验是自然的扩展方向。第四，PC-RCO在极低DVL噪声（$\sigma_v < 0.01$ m/s）下未显示出相对EKF的优势——在这些条件下EKF因其完整的概率框架而表现更优。
